\documentclass[12pt]{article}
\usepackage[utf8]{inputenc}
\usepackage[T1]{fontenc}
\usepackage{hyperref}
\title{Phuza Language Resources for K-12 Curriculum Development}
\author{Compiled by Claude Sonnet 5.5 (claude-sonnet-5-5) for Andr\'as Kornai}
\date{}
\begin{document}
\maketitle

\section{Introduction}
Phuza (ypz) is a Southeastern Loloish language of the Sino-Tibetan family. It is spoken by roughly 6,000 people in southeast Gejiu and southwest Mengzi counties, Yunnan, China. Wikipedia treats it as one of four principal varieties of the Phupha (Downriver Phula) cluster, which may be considered distinct languages. Very little published material exists.

\section{Primary Educational Resources}
No primers, textbooks or school programmes for Phuza turned up in the searches. Wikipedia states that Phuza is not taught in school. It notes that only Alugu, another Phula variety, is taught in primary schools.

\section{Linguistic and Reference Materials}
\begin{itemize}
\item Wikipedia on the Phupha language cluster: \url{https://en.wikipedia.org/wiki/Phupha_language}
\item Wikipedia on the Southeastern Loloish languages: \url{https://en.wikipedia.org/wiki/Southeastern_Loloish_languages}
\item Max Planck numeral page for Phuza, with data collected by Jamin Pelkey in 2008 from a Bujibai (Lengquan, Mengzi) dialect: \url{https://lingweb.eva.mpg.de/channumerals/Phuza.htm}
\end{itemize}

\section{Religious and Community Resources}
\begin{itemize}
\item Joshua Project profile of the Phuza in China (population 6,500): \url{https://joshuaproject.net/people_groups/20612}
\item Joshua Project language page for ypz: \url{https://joshuaproject.net/index.php/languages/ypz}
\end{itemize}

\section{Digital and Audio Resources}
None were found in the searches.

\section{Curriculum Development Context}
Phuza has no known school presence and almost no published documentation. Any curriculum would first need community-led documentation and orthography work. Pelkey (2011) is cited as a study of a Phuza dialect, but no link to it appeared in the search results.
\end{document}